\subsection*{Table S4A. Cross-species integration characteristics.}\phantomsection\label{tab:S4A}
\noindent Species and condition ASW were independently calculated on the same 10,000-observation subsample of the joint 30-dimensional latent space. Their values (0.0346643 and 0.0346610, respectively) both round to 0.0347; the matching displayed values are not duplicate measurements. Low ASW describes label mixing and does not alone establish preservation of biological condition structure. The validation AUROC of 0.9219 is for held-out Walesky mouse cells in the cross-species coordinate system, rather than human donors. Cell/\allowbreak{}nucleus and retained gene-symbol-pair counts are integers; ASW and AUROC are dimensionless.

\begingroup\small\renewcommand{\arraystretch}{1.25}
\noindent\begin{tabularx}{\textwidth}{>{\raggedright\arraybackslash}p{0.38\textwidth}>{\raggedright\arraybackslash}p{0.12\textwidth}X}
\toprule Metric & Value & Note \\ \midrule
Retained unique mouse--human gene-symbol pairs & 13,122 & Mappings present in both Walesky and Matchett expression matrices \\
Walesky mouse hepatocytes & 6,582 & Mouse training/\allowbreak{}reference cells \\
Matchett human hepatocyte nuclei & 49,984 & 10 APAP-ALF donors + 9 healthy donors \\
Joint cross-species cells/\allowbreak{}nuclei & 56,566 & Combined cross-species scVI dataset \\
ASW (species labels) & 0.0347 & Computed in 30-dimensional cross-species scVI latent space \\
ASW (condition labels) & 0.0347 & Cross-species latent-space diagnostic \\
Cross-species TrajRisk validation AUROC & 0.9219 & Held-out Walesky validation set in cross-species coordinate system \\
\bottomrule\end{tabularx}\endgroup\par\medskip
